\origpage{40}{ROBERTSON AND SEYMOUR}
\noindent of members of $\mathscr{P}$. Thus, $V(\mathscr{P}) \subseteq \bar{\Omega}$. We say that $\mathscr{P}$ is \textit{crooked} if for every $P \in \mathscr{P}$ with ends $u, v$ say, there exist $u', v' \in V(\mathscr{P})$ such that $u, u', v, v'$ are distinct and occur in that order in $\Omega$. A \textit{cross} is a crooked transaction of cardinality 2.

Let $\alpha$ be an arrangement in a surface $\Sigma$ of segregation $\mathscr{S}$ of $G$, and let $F \subseteq \Sigma$ be an $\alpha$-normal $O$-arc. Let $X = \{v \in V(\mathscr{S}) : \alpha(v) \in F\}$ and $X' = \{\alpha(v) : v \in X\}$. Let $\Omega'$ be a natural order of $X'$ from $F$, and let $\Omega$ be the cyclic permutation with $\bar{\Omega} = X$ mapped to $\Omega'$ by $\alpha$. We write $\Omega = \alpha^{-1}(F)$, for convenience. The proof of [3, theorem (11.11)] may easily be adapted to yield the following.

\stmt
(8.1) \textit{A society $(G, \Omega)$ has no cross if and only if there is a 3-segregation $\mathscr{S}$ of $G$ with $\bar{\Omega} \subseteq V(\mathscr{S})$ and an arrangement $\alpha$ of $\mathscr{S}$ in a closed disc $\Delta$ with $\Omega = \alpha^{-1}(bd(\Delta))$.}

A \textit{neighbourhood} is a triple $(G, \Omega_1, \Omega_2)$ where $G$ is a graph and $\Omega_1, \Omega_2$ are cyclic permutations in $V(G)$. If $(G, \Omega_1)$ is a society and $(G_1, G_2)$ is a separation of $G$ with $\bar{\Omega}_1 \subseteq V(G_1)$, and $\Omega_2$ is a cyclic permutation with $\bar{\Omega}_2 = V(G_1 \cap G_2)$ , we say that $(G, \Omega_1)$ is the \textit{composition} of $(G_2, \Omega_2)$ and $(G_1, \Omega_1, \Omega_2)$. A neighbourhood $(H, \Omega_1, \Omega_2)$ is \textit{rural} if $H$ is a drawing in a closed cylinder $\Sigma$ with ends $F_1, F_2$ and $\Omega_i$ is a natural order of $V(H) \cap F_i$ from $F_i$ for $i = 1, 2$. (The \textit{ends} of a closed cylinder $\Sigma$ are the two $O$-arcs in $bd(\Sigma)$.) This definition of ``rural'' is slightly different from that in [3], where $\Sigma$ was not required to be a manifold.

Let $(H, \Omega_1, \Omega_2)$ be a rural neighbourhood, in a closed cylinder $\Sigma$ with ends $F_1, F_2$. It has \textit{breadth} $\ge q$ if there are $q$ circuits of $H$, mutually disjoint, and each winding around $\Sigma$, in the natural sense. The following is [3, theorem (10.4)].

\stmt
(8.2) \textit{Let $\pi \ge 0$ be an integer with $\pi \ne 1, 3$, and let $(G, \Omega_1)$ be the composition of a society $(G_2, \Omega_2)$ with a rural neighbourhood $(G_1, \Omega_1, \Omega_2)$ of breadth $\ge \pi + 5$. Suppose that $X \subseteq \bar{\Omega}_1$ and there are $2\pi + 8$ mutually disjoint paths of $G_1$, each from $X$ to $\bar{\Omega}_2$ with no vertex in $\bar{\Omega}_1$ except its initial vertex. If there is a crooked transaction $\mathscr{P}$ in $(G, \Omega_1)$ with cardinality $\ge \pi$, then $\mathscr{P}$ may be chosen so that in addition $V(\mathscr{P}) \subseteq X$ and no member of $\mathscr{P}$ has any internal vertex in $\bar{\Omega}_1$.}

We remark that, in the case $\pi = 2$ (which, for (4.1), is all we need of (8.2)), one can improve (8.2) numerically as follows.

\stmt
(8.3) \textit{Let $(G, \Omega_1)$ be the composition of a society $(G_2, \Omega_2)$ with a rural neighbourhood $(G_1, \Omega_1, \Omega_2)$ of breadth $\ge 3$. Suppose that $X \subseteq \bar{\Omega}_1$ and there are $4$ mutually disjoint paths in $G_1$, each from $X$ to $\bar{\Omega}_2$ with no vertex in $\bar{\Omega}_1$ except its initial vertex. If there is a cross $\mathscr{P}$ in $(G, \Omega_1)$ then $\mathscr{P}$ may be}
